%%%%%%%%%%  scan idx 680  |  folio 667  %%%%%%%%%%
\section*{Exposé XIV. Problèmes ouverts en théorie des intersections}
\begin{center}
par A. Grothendieck
\end{center}

Dans le présent exposé nous passons en revue quelques-uns des problèmes liés aux questions rencontrées dans ce séminaire, et qui mériteraient d'être étudiés. Nous laissons de côté les problèmes liés à la \emph{théorie locale} des intersections à la Serre [14], que nous n'avons pas eu à utiliser dans le Séminaire, et renvoyons notamment à loc. cit. pour les deux conjectures fondamentales qui restent ouvertes dans cette direction (loc. cit. Chap V \S~4). Nous passons sous silence aussi les problèmes, fort importants pour les applications arithmétiques, liés à la \emph{structure des groupes de classes de cycles}, sur un schéma algébrique projectif et lisse notamment; nous renvoyons pour la discussion de ces problèmes à S. Kleiman, \emph{Algebraic cycles and the Weil conjectures}, dans \emph{Dix exposés sur la cohomologie des schémas} (Masson-North Holland Publishing Cie). Dans le no. 4 du présent exposé, nous donnons également quelques compléments au Séminaire, concernant notamment l'\emph{anneau de Chow}, qui n'avaient trouvé leur place dans les exposés précédents.

\section*{1. Opérations $\Lambda^i$ dans la catégorie dérivée $D(X)$\footnote{cf. note de la page 23.}}

Soit $X$ un topos localement annelé, et soit $K^\bullet(X)$ l'anneau des classes de complexes parfaits sur $X$ (Exp. IV). On aimerait définir dans $K^\bullet(X)$ une structure de $\lambda$-anneau (Exp. V), fonctorielle en $X$, telle que l'homomorphisme naturel
\[
 K^\circ_{\mathrm{naif}}(X)\longrightarrow K^\bullet(X)
\]
(où le premier membre est l'anneau des classes de modules localement libres sur $X$, muni de la $\lambda$-structure définie dans Exp. V) soit un $\lambda$-homomorphisme. Dans ce but, il y aurait lieu de définir des foncteurs (non addifs en général de ``puissance extérieure'')
\[
 \Lambda^i:D^-(X)\longrightarrow D^-(X),
\]
transformant complexes pseudo-cohérents (resp. parfaits) en complexes pseudo-cohérents (resp. parfaits), et ayant les propriétés formelles nécessaires, vis-à-vis des triangles exacts, pour définir une $\lambda$-structure sur $K^\bullet(X)$ par la formule
\begin{equation}\tag{1.1}
 \lambda^i(\cl(L^\bullet))=\cl(\Lambda^i L^\bullet).
\end{equation}

%%%%%%%%%%  scan idx 681  |  folio 668  %%%%%%%%%%
Lorsque $X$ est de caractéristique nulle, i.e. son faisceau structural est un faisceau de $\Q$-algèbres, le problème se résoud simplement en prenant
\begin{equation}\tag{1.2}
 \Lambda^i L^\bullet=\left((L^\bullet)^{\otimes i}\right)_{\mathrm{alt}},
\end{equation}
où l'on considère la puissance tensorielle $i$-ème (au sens des catégories dérivées) de $L^\bullet$ comme étant un objet de $D^-(\mathcal O_i)$, $\mathcal O_i$ étant l'algèbre du groupe symétrique $\mathfrak S_i$ à coéfficients dans $\mathcal O_X$, et où l'exposant alt désigne le passage à la ``partie alternée sous l'action de $\mathfrak S_i$'' (opération qui a un sens grâce à l'hypothèse de caractéristique nulle). Une formule de ce type était utilisée déjà dans [8] (cité par la suite [RRR]), pour la première démonstration du théorème de Riemann--Roch.

Sans hypothèse de caractéristique, lorsque $L^\bullet$ est nul en degrés $>0$, alors par Dold--Puppe [6] $L^\bullet$ est canoniquement homotope au complexe associé à un module semi-simplicial $L^{\prime\bullet}$ canoniquement associé à $L^\bullet$, et un candidat assez naturel pour $\Lambda^iL^\bullet$ (pour $L^\bullet$ plat, disons) serait le complexe associé au module semi-simplicial $\Lambda^iL^{\prime\bullet}$ (dont la composante de degré $n$ est $\Lambda^i L^{\prime}_n$). Mais pour étendre la définition à des complexes de $D^-(X)$ qui ne sont plus supposés nuls en degrés $>0$, il faudrait étudier le comportement de l'opération précédente par rapport à l'opération de translation $L^\bullet\mapsto L^\bullet[1]$ (ou \emph{suspension}, dans la terminologie de [6]); malheureusement, il ne semble pas qu'on ait un isomorphisme
\[
 \Lambda^i(L^\bullet[1])\simeq \Lambda^i(L^\bullet)[i]
\]
qui permettrait de faire de façon triviale l'extension désirée.

\section*{2. La formule de Riemann--Roch sans hypothèses projectives}

\textbf{2.1.} Une fois en possession d'une $\lambda$-structure naturelle sur les $K^\bullet(X)$, il est possible de formuler le théorème de Riemann--Roch, sous la forme envisagée dans Exp. VIII, pour n'importe quel morphisme $f:X\to Y$ propre et d'intersection complète, (et même, plus généralement, dans le cas d'un schéma relatif $X$ sur un topos localement annelé $Y$, $X$ étant encore soumis aux conditions relatives correspondantes). La classe de Todd de $f$ se définit grâce à la classe dans $K^\bullet(X)$ du complexe cotangent relatif $T_f$ pour $f$, construit dans Exp. 0, 4.4, lequel complexe est ici parfait grâce à l'hypothèse locale faite sur $f$. Il semble clair que la démonstration de la formule de Riemann--Roch dans ce cas général demandera l'introduction d'idées essentiellement nouvelles, et il est possible que ces mêmes idées pourraient donner la clef d'une
%%%%%%%%%%  scan idx 682  |  folio 669  %%%%%%%%%%
démonstration de la formule analogue dans le cas analytique complexe (ou rigide-analytique !). Rappelons que dans le cas analytique complexe, seul le cas d'un faisceau localement libre $F$ sur une variété non singulière compacte $X$, au-dessus d'un $Y$ ponctuel, est actuellement connu (grâce à la formule de l'index de Atiyah--Singer), le cas où $F$ est seulement supposé cohérent semblant échapper encore aux techniques connues jusqu'à présent.

\textbf{2.2.} Comme ingrédient essentiel de la formule de Riemann--Roch, rappelons aussi pour mémoire le \emph{problème de finitude} soulevé dans Exp. III: si $f:X\to Y$ est un morphisme propre, et $L^\bullet$ un complexe sur $X$ qui est pseudo-cohérent relativement à $Y$, alors $Rf_*(L^\bullet)$ est-il pseudo-cohérent? Une démonstration générale de cette conjecture mettra probablement en oeuvre des idées algébriques qui permettront de donner une démonstration du ``théorème'' de finitude de Grauert [7] sous des conditions notablement plus générales que dans loc. cit., en se plaçant d'emblée sur des espaces (ou topos) de base topologiquement annelés en $\C$-algébres $Y$ plus généraux que des espéces analytiques (englobant le cas où $Y$ serait un espace localement compact avec le faisceau des fonctions continues complexes, ou une variété $C^m$ avec le faisceau des fonctions complexes $C^m$, ou un espace analytique réel, etc).

\textbf{2.3.} Mentionnons enfin qu'il y aurait lieu de remplacer dans ces questions l'hypothèse de propreté sur le morphisme $f$ par l'hypothèse que les faisceaux de cohomologie $H^i(L^\bullet)$ des complexes envisagés sur $X$ aient leurs groupes propres sur $Y$.

\section*{3. Formule de Riemann--Roch ``sans dénominateurs'' pour une immersion}

\textbf{3.1.} Soit $f:Y\to X$ une immersion fermée régulière (Exp. VII), de faisceau conormal $\underline N$. Supposons pour simplifier que $X$ admet un faisceau inversible ample. La formule de Riemann--Roch pour $y\in K^\bullet(Y)$
\[
 \ch_X(f_*(y))=f_*\bigl(\ch_Y(y)\operatorname{Todd}(-\check{\underline N})\bigr),
\]
jointe à la formule générale (Exp. VII)
\[
 f_*(y')\,f_*(y'')=f_*\bigl(y'y''c_d(\underline N)\bigr),\qquad y',y''\in \operatorname{Gr}^\bullet(Y)_\Q,
\]
(en supposant $f$ de codimension $d$), permet d'exprimer les classes de Chern de $f_*(y)$ dans $\operatorname{Gr}^\bullet(X)_\Q$ sous la forme
\begin{equation}\tag{3.1}
 c_i(f_*(y))=f_*\bigl(P_i((c_\alpha(y))_\alpha,(c_\beta(\underline N))_\beta)\bigr),
%%%%%%%%%%  scan idx 683  |  folio 670  %%%%%%%%%%
\end{equation}
où $P_i$ est un polynôme universel en les $c_\alpha(y)$, $c_\beta(\underline N)$, polynôme qui a priori est à coefficients rationnels. En fait, on constate [RRR] que ces polynômes sont à coefficients \emph{entiers}. Cela soulève la question de la validité d'une formule de Riemann--Roch ``sans dénominateurs''.

En fait, une telle formule pour $c_i(f_*(y))$ ne peut être cherchée en général dans $\operatorname{Gr}^\bullet(X)$ lui-même, vu que l'homomorphisme $f_*:K^\bullet(Y)\to K^\bullet(X)$ n'est pas compatible lui-même avec les $\lambda$-filtrations modulo décalage par $d$ (4.6), mais l'est seulement à torsion près; de sorte qu'il n'est pas possible de définir un homomorphisme
\[
 f_*:\operatorname{Gr}^\bullet(Y)\longrightarrow \operatorname{Gr}^\bullet(X)
\]
donnant l'homomorphisme naturel $f:\operatorname{Gr}^\bullet(Y)_\Q\to\operatorname{Gr}^\bullet(X)_\Q$, et permettant de donner un sens à la formule (3.1) comme une formule dans $\operatorname{Gr}^\bullet(X)$, et non seulement dans $\operatorname{Gr}^\bullet(X)_\Q$.

Mais lorsque $X$ est supposé régulier, donc $K^\bullet(X)\simeq K_\bullet(X)$ (IV 2.5), on peut munir $K^\bullet(X)$ de la \emph{filtration topologique} par la codimension des supports, héritée de la filtration correspondante de $K_\bullet(X)$. Il est plausible que cette filtration est compatible avec la structure d'anneau, fait qu'on peut prouver tout au moins [RRR] lorsque $X$ est lisse sur un corps $k$, en utilisant la théorie de l'anneau de Chow [4]. Quand il en est ainsi, le gradué associé
\[
 \Gtop^\bullet(X)
\]
est un \emph{anneau gradué}, et supposant également $Y$ régulier, on a fait ce qu'il fallait pour avoir un homomorphisme naturel de degré $d$ de groupes gradués
\[
 f_*=f_*^{\mathrm{top}}:\Gtop^\bullet(Y)\longrightarrow \Gtop^\bullet(X).
\]
Comme la $\lambda$-filtration de $K^\bullet(X)$ est plus fine que la filtration topologique par la codimension qu'on vient de considérer (X 1.3.3), on trouve un homomorphisme canonique de groupes gradués (et même d'anneaux gradués)
\[
 \operatorname{Gr}^\bullet(X)\longrightarrow \Gtop^\bullet(X),
\]
qui est d'ailleurs un isomorphisme mod torsion (Exp. VIII). Ceci permet, comme dans [RRR], de considérer une théorie des classes de Chern sur $X$ à valeurs dans $\Gtop^\bullet(X)$, et de même sur $Y$, et de donner ainsi un sens à la formule (3.1) comme une formule dans $\Gtop^\bullet(X)$. Sous cette forme, cette formule est prouvée dans [RRR] lorsque $X$ est de \emph{caractéristique nulle}, tout au moins si $X$ est lisse sur un corps (hypothèse servant uniquement à assurer que la filtration topologique de $K^\bullet(X)$ est compatible avec la structure d'anneau). L'hypothèse de caractéristique nulle est utilisée
%%%%%%%%%%  scan idx 684  |  folio 671  %%%%%%%%%%
essentiellement pour disposer de formules du type (1.1), (1.2), et la même démonstration essentiellement devrait s'appliquer sans hypothèse de caractéristique, une fois résolue dans un sens affirmatif la question soulevée au no. 1.

\textbf{3.2.} Lorsque $X$ est lisse sur un corps, on dispose aussi de l'anneau de Chow
\[
 A(X),
\]
et d'une théorie des classes de Chern sur $X$ à valeurs dans $A(X)$ [9], plus fine que la théorie à valeurs dans $\Gtop^\bullet(X)$ qu'on vient d'introduire, ces dernières classes de Chern étant images des classes de Chern style Chow par l'homomorphisme canonique surjectif (cf. 4.1)
\[
 \varphi:A(X)\longrightarrow \Gtop^\bullet(X).
\]
Cet homomorphisme est un isomorphisme mod. torsion, mais pas un isomorphisme en général (4.7). Ceci dit, si $X$ et $Y$ sont lisses, l'inclusion $f:Y\to X$ définit encore un homomorphisme de degré $d$ de groupes gradués
\[
 f_*:A(Y)\longrightarrow A(X),
\]
ce qui permet de donner un sens à la formule (3.1) comme une formule dans $A^i(X)$\footnote{La formule, ainsi que la variante signalée dans 3.3 a été prouvée depuis par P. Jouanolou (nov.\ 1969). Cf.}. Cette formule se réduit d'ailleurs à $0=0$ pour $i<d$, et ne présente donc d'intérêt que pour $i\geq d$. Même pour $i=d=2$ cette formule n'est pas prouvée à l'heure actuelle (cf. 4.3), ni la formule moins précise qui s'en déduit en remplaçant l'équivalence rationelle des cycles algébriques par l'équivalence algébrique. Contrairement à la formule analogue dans $\Gtop^\bullet(X)$, il ne semble pas y avoir de raison bien convainquante pour supposer que la formule sans dénominateurs doive être vraie dans $A(X)$. Une bonne façon de la tester, vue notre ignorance quasi totale sur la nature des groupes $A^i(X)$ pour $i\geq 2$, semble être de tester les formules qu'on peut en déduire dans des $A^1(Z)=\Pic(Z)$, par exemple en transformant les deux membres de (3.1) par
\[
 g_*:A(X)\longrightarrow A(Z),
\]
où $g:X\to Z$ est un morphisme propre dans un $Z$ quasi-projectif lisse de dimension $\dim X-i+1$.

\textbf{3.3.} Enfin, pour donner un sens à (3.1), on peut également utiliser la théorie cohomologique des classes de Chern (SGA 5 VII), chaque fois qu'on dispose d'\emph{homomorphismes de Gysin}
\[
 f_*:H^p(Y,\mu_n^{\otimes i})\longrightarrow H^{p+2d}(X,\mu_n^{\otimes(i+d)}).
\]
(cf. no.~6), ce qui est le cas tout au moins si $Y$ et $X$ sont tous deux lisses sur un schéma $S$ (grâce au théorème de pureté cohomologique relatif (SGA 4 XVI 3.7), et,
%%%%%%%%%%  scan idx 685  |  folio 672  %%%%%%%%%%
conjecturalement, si $Y$ et $X$ sont tous deux réguliers (d'après le théorème de pureté cohomologique absolu, démontré tout au moins pour $X$ excellent de caractéristique nulle (SGA 4 XIX 3.2)). Lorsque $X$ et $Y$ sont lisses sur le corps des complexes, on peut également utiliser la classe de Chern entière sur l'espace localement compact $X(\C)$ des points complexes de $X$. La formule (3.1) est établie dans ce cas par Atiyah et Hirzebruch [1], qui se placent plus généralement dans le cas d'une immersion de variétés analytiques complexes (en travaillant avec le $K^\circ$ ``naif'') ; il en résulte que la formule (3.1) $\ell$-adique est valable chaque fois que $Y$ et $X$ sont lisses sur un corps \emph{algébriquement clos et de caractéristique nulle}, en utilisant le ``principe de Lefschetz'' et les résultats généraux de SGA 4. Cela rend donc la validité de la formule (3.1) cohomologique extrêmement plausible; mais on notera qu'elle n'est pas démontrée à l'heure actuelle, même si $X$ et $Y$ sont lisses sur un corps $k$ de caractéristique nulle $k$, lorsque $k$ n'est pas algébriquement clos.

\textbf{3.4.} On peut également travailler avec la théorie cohomologique des classes de Chern, style De Rham ou style Hodge [13]. Le problème de la validité de (3.1) dans ce cas ne se pose que si $X$ n'est pas de caractéristique nulle, puisque l'on sait de toutes façons (no.~6) que la formule en question est vraie modulo torsion. Pour les schémas lisses quasi-projectifs sur un corps $k$, la formule (3.1) version De Rham ou Hodge serait d'ailleurs conséquence de la formule (3.1) version classes de cycles mod. équivalence algébrique, de sorte que la version cohomologique De Rham ou Hodge donne une autre façon abordable de tester la validité de (3.1) version ``classes de cycles''.

\textbf{3.5.} La démonstration [1] de la version cohomologique transcendante de (3.1) est de nature essentiellement transcendante, en faisant usage de voisinages tubulaires de $Y$ dans $X$. Cela suggère que la démonstration d'une formule du type (3.1) pourrait être liée à la considération systématique de groupes $K^\bullet$ et de groupes de cohomologie \emph{à support dans} $Y$, considération qui pourrait être un moyen technique suffisant de ``localisation autour de $Y$''.

\section*{4. Relations entre $K^\bullet(X)$ et l'anneau de Chow $A(X)$}

\textbf{4.1.} Supposons pour simplifier que $X$ soit un schéma lisse et quasi-projectif sur un corps $k$ (cf. no.~8 pour un cas plus général), de sorte qu'on
%%%%%%%%%%  scan idx 686  |  folio 673  %%%%%%%%%%
dispose de l'anneau de Chow $A(X)$ des classes de cycles à équivalence rationelle près [4]. Comme variantes de cet anneau, défini d'ailleurs sous des considérations sensiblement plus générales, on a les anneaux $\operatorname{Gr}^\bullet(X)$ (gradué associé à la $\lambda$-filtration de $K^\bullet(X)$) et $\Gtop^\bullet(X)$ (gradué associé à la filtration de $K^\bullet(X)\simeq K_\bullet(X)$ par la codimension des supports). On a des théories des classes de Chern à valeurs dans chacun de ces trois anneaux (cf. [9] pour le cas de $A(X)$, et ce séminaire pour $\operatorname{Gr}^\bullet(X)$, qui s'envoie dans $\Gtop^\bullet(X)$). On a également une théorie cohomologique ($\ell$-adique, disons, où $\ell$ est premier à la caractéristique de $k$) des classes de Chern (SGA 5 VII)~:
\[
 c_i(E)\in H^{2i}(X,\Zell(i)),
\]
(où $E$ est un fibré vectoriel sur $X$, ou encore un élément quelconque de $K^\bullet(X)$). Ces théories cohomologiques sont reliées à l'aide des homomorphismes canoniques (compatibles avec les structures multiplicatives)
\begin{equation}\tag{4.1}
\begin{tikzcd}[column sep=large,row sep=large]
 H^{2\bullet}(X,\Zell(\bullet)) & {} \\
 A(X) \arrow[u,"\gamma"] \arrow[r,"\varphi"] & \Gtop^\bullet(X) \\
 {} & \operatorname{Gr}^\bullet(X) \arrow[u,"j"']
\end{tikzcd}
\end{equation}
où on a posé
\[
 H^{2\bullet}(X,\Zell(\bullet))=\bigoplus_i H^{2i}(X,\Zell(i)).
\]
L'homomorphisme $j$ provient du fait que la $\lambda$-filtration de $K^\bullet(X)$ est plus fine que la filtration topologique par la codimension du support (Exp. X), et c'est un isomorphisme modulo torsion (Exp. VII 4.11). L'homomorphisme $\varphi$ est défini par exemple dans [RRR], et associe à la classe du cycle de codimension $i$ associé à un sous-schéma fermé $Y$ de $X$ codimension $i$, l'élément $\cl^i(\mathcal O_Y)$ de
\[
 \Fil^i_{\mathrm{top}}K^\bullet(X)/\Fil^{i+1}_{\mathrm{top}}K^\bullet(X).
\]
C'est un homomorphisme surjectif, dont nous allons voir encore (4.2) que c'est un isomorphisme modulo torsion. Ainsi, \emph{modulo torsion les trois anneaux} $A(X)$, $\Gtop^\bullet(X)$ \emph{et} $\operatorname{Gr}^\bullet(X)$ \emph{sont canoniquement isomorphes}, donc \emph{à posteriori} les formules de Riemann--Roch exprimées à l'aide de l'un ou l'autre de ses anneaux (tensorisés par $\Q$) sont équivalentes. L'anneau gradué $\operatorname{Gr}^\bullet(X)$, surtout utilisé dans ce séminaire comme substitut à l'anneau de Chow, a l'avantage appréciable d'être défini pour tout topos localement annelé, et en particulier pour tout schéma (
%%%%%%%%%%  scan idx 687  |  folio 674  %%%%%%%%%%
régulier ou non). Il a par contre le désavantage (4.6) sur $A(X)$, $\Gtop^\bullet(X)$ de ne pas avoir tel quel un caractère covariant par rapport aux morphismes propres (même pour des schémas lisses quasi-projectifs sur un corps)~; il n'acquiert ce caractère qu'une fois tensorisé par $\Q$.

Autre avantage de $A(X)$ sur ses concurrents $\Gtop^\bullet(X)$ et $\operatorname{Gr}^\bullet(X)$, c'est qu'on dispose d'un homomorphisme de $A(X)$ dans $H^{2\bullet}(X,\Zell(\bullet))$, permettant d'interpréter les classes de Chern cohomologiques comme images des classes de Chern à valeurs dans $A(X)$. Cet homomorphisme en général ne se factorise pas par $\Gtop^\bullet(X)$ (4.7), et il est pour le moins douteux qu'il existe en général un homomorphisme de $\operatorname{Gr}^\bullet(X)$ dans $H^{2\bullet}(X,\Zell(\bullet))$ qui transforme classes de Chern en classes de Chern (Cf. no 5). Bien entendu, ces difficultés disparaissent lorsqu'on tensorise par $\Q$, et on trouve un homomorphisme
\begin{equation}\tag{4.2}
 \Gtop^\bullet(X)\longrightarrow H^{2\bullet}(X,\Qell(\bullet))=
 H^{2\bullet}(X,\Zell(\bullet))\otimes_{\Zell}\Qell
\end{equation}
qui transforme classes de Chern en classes de Chern.

\textbf{4.2.} Pour établir que l'homomorphisme surjectif $\varphi$ de (4.1) est un isomorphisme modulo torsion, on utilise l'homomorphisme de Chern (déduit de la théorie des classes de Chern [9])
\[
 K^\bullet(X)\longrightarrow \Ch(A(X)),
\]
(où le deuxième membre est l'anneau de Chern de $A(X)$ défini dans Exp.~V), qui est compatible comme on voit aisément avec la filtration topologique de $K^\bullet(X)$ et la filtration naturelle de $\Ch(A(X))$, et donne par passage aux gradués associés un homomorphisme de groupes gradués
\begin{equation}\tag{4.3}
 \psi:\Gtop^\bullet(X)\longrightarrow A(X).
\end{equation}
Ceci posé, on va prouver les formules
\begin{equation}\tag{4.4}
\left\{
\begin{aligned}
 \psi\varphi(x)&=(-1)^{i-1}(i-1)!\,x\quad \bmod \text{ torsion}\quad \text{pour } x\in A^i(X),\\
 \varphi\psi(x)&=(-1)^{i-1}(i-1)!\,x\quad \bmod \text{ torsion}\quad \text{pour } x\in \Gtop^i(X);
\end{aligned}
\right.
\end{equation}
dont la première implique d'ailleurs la seconde, grâce au fait que $\varphi$ est surjectif. Ces formules impliquent évidemment que $\varphi$ et $\psi$ sont des isomorphismes modulo torsion, et qu'à un facteur $(-1)^{i-1}(i-1)!$ près en degré $i$, ils sont inverses l'un de l'autre.

Pour prouver la première formule (4.4), on peut évidemment supposer que $x$ est défini par un sous-schéma fermé irréductible $Y$ de $X$ de codimension $i$, auquel cas la formule équivaut à
\begin{equation}\tag{4.5}
 c_i(\mathcal O_Y)=(-1)^{i-1}(i-1)!\,\cl(Y)\quad \bmod
%%%%%%%%%%  scan idx 688  |  folio 675  %%%%%%%%%%
\text{ torsion},
\end{equation}
où dans le premier membre $\mathcal O_Y$ est considéré comme un faisceau cohérent sur $X$, définissant donc un élément de $K^\bullet(X)$, et dans le deuxième membre $\cl(Y)$ désigne la classe de $Y$ dans $A^i(X)$. Comme $A^i(X)$ ne change pas en enlevant de $X$ une partie fermée de codimension $>i$, on est ramené au cas où $Y$ est régulier, donc lisse sur le corps de base $k$ si ce dernier est parfait. Mais dans ce dernier cas, la formule (4.5) n'est autre que la formule de Riemann--Roch pour l'inclusion $Y\to X$ et l'élément $1$ de $K^\bullet(Y)$, formule écrite en utilisant la théorie des classes de Chern à valeurs dans l'anneau de Chow $A(X)$. Cette formule est établie dans [2] sous les conditions envisagées ici, par une méthode très voisine de celle utilisée dans ce séminaire pour la formule analogue utilisant $\operatorname{Gr}^\bullet(X)$. Cela établit donc, du moins pour $k$ parfait, la formule (4.4) et par suite le fait que $\varphi$ soit un isomorphisme mod. torsion (qui entraîne à son tour, à posteriori, que la formule de Riemann--Roch ``avec dénominateurs'', pour un morphisme propre de schémas algébriques lisses sur $k$, est essentiellement la même, qu'on l'écrive en utilisant $A(X)$ comme dans [2], ou en utilisant $\operatorname{Gr}^\bullet(X)$ comme dans ce séminaire !). Le cas où on ne fait pas d'hypothèse sur $k$ se ramène aisément au cas où $k$ est algébriquement clos, en utilisant le fait que pour une extension finie $k'$ de degré $n$ de $k$, le noyau de l'application naturelle $A(X)\xrightarrow{f_X^*}A(X\otimes_k k')$ est annulé par $n$ (comme on voit en utilisant l'homomorphisme trace $f_{X*}:A(X\otimes_k k')\to A(X)$, satisfaisant $f_{X*}f_X^*=n\,\id$).\footnote{Elle est donc établie grâce à P. Jouanolou pour $X$ lisse quasi-projectif sur un corps (cf. note p.5).}

\textbf{4.3.} Il semble assez plausible que les formules (4.4) soient vraies telles quelles et non seulement modulo torsion, ou ce qui revient au même, que la formule (4.5) soit vraie non seulement modulo torsion. Cette question est un cas très particulier de celle de la validité dans $A(X)$ de la ``formule de Riemann--Roch sans dénominateurs'' (3.1), pour une immersion $Y\to X$, (savoir le cas où $i=d$). On peut considérer la validité de ce cas particulier comme assez plausible. En effet, nous avons signalé déjà au no. 3 que lorsque $k$ est de caractéristique nulle, la formule de Riemann--Roch sans dénominateurs est valable tout au moins dans $\Gr^{\bullet}_{\mathrm{top}}(X)$, ce qui implique dans ce cas que les deux membres de (4.5) ont même image dans $\Gr^{\bullet}_{\mathrm{top}}(X)$, ou, ce qui revient au même, que la deuxième des formules (4.4) est valable; ce résultat s'étendra de lui-même au cas de la caractéristique quelconque, une fois résolu de
%%%%%%%%%%  scan idx 689  |  folio 676  %%%%%%%%%%
façon satisfaisante la question soulevée au no 1. Comme support pour la formule (4.5) dans $A(X)$ lui-même, signalons le fait qu'il est clair à priori que la restriction du premier membre à $X-Y$ est nulle, donc, $Y$ étant supposé irréductible, le premier membre doit être à priori un multiple entier de $\cl(Y)$:
\[
 c_i(Y)=n(Y)\cl(Y).
\]
Cela prouve déjà la formule exacte (4.5) lorsque $\cl(Y)$ n'est pas un élément de torsion de $A^i(X)$. Il est immédiat que tel est le cas si $X$ est projectif (et non seulement quasi-projectif) sur le corps $k$, en notant que si $\xi\in A^1(X)$ est une polarisation de $X$, alors $\deg(Y\cdot\xi^i)>0$; cela prouve la formule voulue dans ce cas. Il en résulte que cette formule est encore valable chaque fois que $X$ est isomorphe à un sous-schéma ouvert d'un schéma projectif lisse sur $k$, ce qui sera le cas lorsqu'on dispose d'une solution convenable du problème de résolution des singularités. Il en est ainsi en particulier si $k$ est de caractéristique nulle, grâce à Hironaka [11].

La validité de la formule exacte (4.5) donc de (4.4) impliquerait que le noyau de
\[
 \varphi^i:A^i(X)\longrightarrow \Gtop^i(X)
\]
est annulé par $(i-1)!$, et en particulier \emph{que $\varphi^2:A^2(X)\to\Gtop^2(X)$ est un isomorphisme}. (Que $\varphi^1$ soit un isomorphisme résulte déjà immédiatement de X 5.3.2). Sauf erreur, le fait que $\varphi^2$ soit un isomorphisme n'est pas prouvé à l'heure actuelle, même dans le cas particulier où $\dim X=3$. C'est hâtivement en tous cas que le rédacteur avait affirmé la validité de (4.4) et (4.5) dans [9, no 4, 3o].

\textbf{4.4.} A l'heure actuelle, le principal avantage de $\Gtop^\bullet(X)$ (ou de $\operatorname{Gr}^\bullet(X)$) sur l'anneau de Chow $A(X)$ est le fait qu'il se prête mieux aux calculs explicites de certaines formules d'intersection entre cycles, dans le cas d'intersections ``excédentaires'' i.e. de dimension trop grande. La définition de la classe d'intersections, à l'aide du lemme de Chow permettant de bouger suffisamment un cycle dans sa classe, se prête en effet malaisément au calcul. Un exemple frappant à cet égard est la formule fondamentale suivante [Exp. VII 2.7], pour une immersion fermée $f:Y\to X$ de schémas quasi-projectifs lisses sur $k$:
\begin{equation}\tag{4.6}
 f^*f_*(y)=y\,c_d(N),
\end{equation}
formule valable dans $\Gtop^\bullet(X)$, $y$ étant dans $\Gtop^\bullet(Y)$, $N$ le faisceau conormal de $Y$
%%%%%%%%%%  scan idx 690  |  folio 677  %%%%%%%%%%
dans $X$, supposé partout de rang $d$. Dans loc. cit., cette formule résulte d'un calcul très simple de faisceaux $\underline{\Tor}$. Grâce à ce qu'on a vu dans 4.2, la formule précédente implique la même formule dans $A(Y)$, mais seulement modulo torsion. Il est assez extraordinaire que la formule (4.6) dans $A(Y)$ lui-même ne soit pas prouvée à l'heure actuelle\footnote{cf. note de la page 23.}, ni même le cas particulier obtenu en faisant $y=1$, donnant la self-intersection de $Y$:
\begin{equation}\tag{4.7}
 f^*(\cl(Y))=c_d(N),
\end{equation}
même si $k$ est le corps des complexes ! On se heurte à un problème analogue lorsqu'on tente de calculer l'anneau de Chow de la variété éclatée $X'$ de $X$ le long de $Y$, en calquant la méthode suivie pour le calcul de $\operatorname{Gr}^\bullet(X')_\Q$ [Exp. VII 4.8], et qui marche également (sans tensoriser par $\Q$) pour le calcul de $\Gtop^\bullet(X')$. La formule-clef pour cette détermination (avec les notations de loc. cit.) est
\begin{equation}\tag{4.8}
 f^*(i_*(y))=j_*\bigl(g^*(y)c_{d-1}(\check{\underline F})\bigr),
\end{equation}
formule établie à l'heure actuelle dans $\Gtop^\bullet(X')$ (pour $y\in\Gtop^\bullet(Y)$) grâce à un calcul de $\underline{\Tor}$ (Exp. VII 3.4), et dont la validité dans $A(X')$ (pour $y\in A(Y)$) resterait à établir.

Signalons que même les analogues $\ell$-adiques des formules (4.6), (4.7), (4.8) ne sont pas prouvés à l'heure actuelle; lorsque $y$ est la classe de cohomologie d'un cycle algébrique, (4.6) et (4.8) sont vraies en tous cas mod. torsion.

\textbf{4.5.} Dans cette section et les suivantes, nous allons donner des exemples où dans le diagramme (4.1) l'homomorphisme $j$ respectivement l'homomorphisme $\varphi$ n'est pas un isomorphisme. Il revient évidemment au même, pour $X$ donné, de dire que
\[
 j:\operatorname{Gr}^\bullet(X)\longrightarrow \Gtop^\bullet(X)
\]
est un isomorphisme, ou un monomorphisme, ou un épimorphisme. D'ailleurs l'image de $j$ est le sous-anneau de $\Gtop^\bullet(X)$ engendré par les classes de Chern des faisceaux localement libres sur $X$. Dans le cas où $K^\bullet(X)$ est engendré comme algébre sur $\Z$ par les classes des faisceaux inversibles sur $X$, on voit tout de suite que le sous-anneau envisagé de $\Gtop^\bullet(X)$ est aussi le sous-anneau engendré par $\Gtop^1(X)$, la composante de degré $1$ de $\Gtop^\bullet(X)$.

Or
%%%%%%%%%%  scan idx 691  |  folio 678  %%%%%%%%%%
prenons pour $X$ une \emph{variété de drapeaux}
\[
 X=G/B,
\]
où $G$ est un groupe algébrique semi-simple simplement connexe sur le corps algébriquement clos $k$, et $B$ un sous-groupe de Borel [5]. Lorsque $k=\C$, alors la décomposition cellulaire de Bruhat de $X$ [5, Exp. 13] implique que l'homomorphisme canonique
\[
 (*)\qquad K^\bullet(X)\longrightarrow K(X(\C))
\]
est un isomorphisme, de même d'ailleurs que l'homomorphisme canonique
\[
 (**)\qquad A(X)\longrightarrow H^\bullet(X(\C),\Z),
\]
où $X(\C)$ est l'espace compact des points complexes de $X$, muni de sa topologie habituelle déduite de la topologie de $\C$. D'autre part, il a été prouvé par L. Hodgkin [12] que $K(X(\C))$ est bien engendré par les classes des faisceaux inversibles sur $X$. Il en est donc de même de $K^\bullet(X)$ lui-même, et il n'est d'ailleurs pas difficile d'en déduire, par une technique de spécialisation utilisant les groupes semi-simples déployés sur $\Z$ (SGA 3), que le même résultat reste valable pour tout corps $k$ algébriquement clos. D'ailleurs, la décomposition cellulaire de $X$ (ou la bijectivité de $(**)$ dans le cas $k=\C$) implique aisément que $A(X)$ est sans torsion, donc grâce à 4.2 l'homomorphisme
\[
 \varphi:A(X)\longrightarrow\Gtop^\bullet(X)
\]
est un isomorphisme. De ces faits il résulte que $j$ est un isomorphisme si et seulement si $A(X)$ est engendré par ses éléments de degré $1$, ou encore, si $k=\C$, si et seulement si la cohomologie entière de $X(\C)$ est engendré par les éléments de degré $2$, i.e. si le classifiant $B_{G(\C)}$ est sans torsion. Il est connu d'autre part [15, p. 217 th. B (1)$\Leftrightarrow$(2), et 2.5] que ceci est le cas si et seulement si $G$ est "sans torsion", ou encore si $G$ est isomorphe à un produit de groupes $\operatorname{Sl}(n)$ et $\operatorname{Sp}(m)$. Il suffit donc de prendre pour $G$ un groupe $\operatorname{SO}(n)$ avec $n\geq 7$, ou un groupe simple exceptionnel, pour avoir un exemple où $j$ n'est pas un isomorphisme.

\textbf{4.6.} On notera que dans cet exemple, il est hors de doute que le "bon" anneau est l'anneau sans torsion
\[
 A(X)\simeq\Gtop^\bullet(X),
\]
ayant comme base les cellules de la décomposition de Bruhat, anneau également isomorphe à la cohomologie entière de $X(\C)$ si $k=\C$; par contre, l'anneau $\operatorname{Gr}^\bullet(X)$ apparait comme un anneau ayant de la torsion de nature assez pathologique. Cet exemple confirme donc la conception suivant laquelle l'invariant $\operatorname{Gr}^\bullet(X)$ n'est guère
%%%%%%%%%%  scan idx 692  |  folio 679  %%%%%%%%%%
intéressant que modulo torsion, i.e. tensorisé par $\Q$. On voit en même temps, à l'aide de cet exemple, que ce n'est que tensorisé par $\Q$ que $\operatorname{Gr}^\bullet(X)$ à un caractère covariant par rapport aux morphismes propres, même en se bornant aux immersions fermées de schémas quasi-projectifs et lisses. En effet, si $X$ est tel que $j$ ne soit pas un isomorphisme, il existe dans $X$ un sous-schéma fermé intègre $Y$ tel que $\cl^i(Y)\in K^\bullet(X)$ ne soit pas dans $\Fil^i(X)$, où $i=\operatorname{codim}(Y,X)$. Nous pouvons prendre $Y$ de dimension minimale, de sorte que $\Fil^{i+1}(X)=\Fil_{\mathrm{top}}^{i+1}(X)$. Si $Z$ est l'ensemble singulier de $Y$, et $X'=X-Z$, $Y'=Y-Z$, on en conclut que $\cl(Y')\in K^\bullet(X')$ n'est pas dans $\Fil^i(X')$ (compte tenu que $K^\bullet(X)\to K^\bullet(X')$ est surjectif et compatible avec les augmentations, et que son noyau est $\operatorname{Im}\bigl(K(Z)\longrightarrow K(X)\bigr)\subset\Fil_{\mathrm{top}}^{i+1}(X)=\Fil_{\mathrm{top}}^{i+1}(X)$). Considérons alors l'inclusion
\[
 f:Y'\longrightarrow X',
\]
qui est une immersion fermée de codimension $i$ de schémas quasi-projectifs et lisses sur $k$, \emph{telle que}
\[
 f_*:K^\bullet(Y')\longrightarrow K^\bullet(X')
\]
\emph{ne soit pas compatible avec les $\lambda$-filtrations (modulo décalage par $i$)}, puisque $f_*(1)$ n'est pas dans $\Fil^i(X')$.

\textbf{4.7.} Gardant les notations de 4.5, notons maintenant qu'il résulte de (Exp. IX) que le théorème de Hodgkin cité équivaut au fait que l'on a un isomorphisme
\[
 K^\bullet(G)\simeq\Z,\qquad\text{donc}\qquad \Gtop^\bullet(G)\simeq\Z,
\]
l'homomorphisme $K^\bullet(G/B)\to K^\bullet(G)$ étant à priori surjectif et ayant comme noyau l'idéal engendré par les $\cl'(L)-1$, où $L$ parcourt les Modules inversibles sur $X=G/B$ associés aux représentations $B\to\Gm$ (qui donnent en fait tous les modules inversibles sur $X$ [5, Exp. 15]). On voit d'autre part par un raisonnement analogue que
\[
 A(G/B)\longrightarrow A(G)
\]
est également surjectif, et que son noyau est l'idéal engendré par les éléments de $A^1(G/B)$ définis par les représentations $B\to\Gm$, donc en fait par $A^1(G/B)$ tout entier.\footnote{Cf. IX 3.4, 2.2.} Cela montre que $A(G)$ est réduit à $A^0(X)=\Z$ si et
%%%%%%%%%%  scan idx 693  |  folio 680  %%%%%%%%%%
seulement si "$G$ est sans torsion", i.e. $G$ est isomorphe à un produit de groupes $\operatorname{Sl}(n)$ et $\operatorname{Sp}(m)$. Donc dans ce cas, et dans ce cas seulement, l'homomorphisme
\[
 \varphi:A(G)\longrightarrow\Gtop^\bullet(G)
\]
est un isomorphisme. Donc si $G$ a de la torsion, le schéma sous-jacent fournit un exemple d'un $\varphi$ non bijectif.

\textbf{4.8.} On trouve de même un exemple où l'homomorphisme $\gamma$ de (4.1)
\begin{equation}\tag{4.9}
 \gamma:A(G)\longrightarrow H^{2\bullet}(G,\Zell(\bullet))
\end{equation}
ne se factorise pas par $\Gtop^\bullet(G)$, ou ce qui revient au même ici, un exemple où cet homomorphisme n'est pas nul en degrés $>0$. Se limitant (comme il est loisible grâce à des théorèmes de spécialisation convenables) au cas où $k=\C$, on est ramené à étudier l'homomorphisme
\begin{equation}\tag{4.10}
 A(G)\longrightarrow H^\bullet(G(\C),\Z),
\end{equation}
déduit d'ailleurs par passage au quotient de l'homomorphisme canonique
\[
 A(G/B)\simeq H^\bullet(G/B(\C),\Z)\longrightarrow H^\bullet(G(\C),\Z).
\]
Il est tentant de conjecturer que l'homomorphisme (4.10) est en fait injectif, puisqu'on sait par ailleurs que les nombres premiers de torsion pour $G(\C)$ sont exactement ceux qui divisent l'ordre du groupe fini $A^+(G)=\coprod_{i\geq0}A^i(G)$. S'il en était ainsi, dès que $G$ aurait de la torsion, l'homomorphisme (4.10) ne se factoriserait pas par $\varphi$, donc pour $\ell$ convenable (de façon précise, pour $\ell$ divisant l'ordre de $A^+(G)$, i.e. $\ell$ coéfficient de torsion de $G$) il en serait de même pour (4.9). On peut du moins affirmer que dans le cas du groupe $G=G_2$ sur le corps $\C$, l'homomorphisme (4.10) n'est pas nul en degré $>0$; de façon précise il existe un élément de $A^3(G)$ dont l'image dans $H^6(G(\C),\Z)$ est d'ordre $2$; il en résulte donc que l'homomorphisme (4.9) ne se factorise pas par $\varphi$ lorsque $\ell=2$. Le fait utilisé concernant $G_2$, qui revient à dire qu'il existe dans $H^6(G/B(\C),\Z)$ un élément dont l'image dans $H^6(G(\C),\Z)$ est non nulle, se trouve dans [3], modulo le dictionnaire connu entre le langage groupes algébriques semi-simples complexes, et groupes de Lie compacts; il m'a été communiqué obligeamment par J.-P. Serre.

%%%%%%%%%%  scan idx 694  |  folio 681  %%%%%%%%%%
\section*{5. Relations entre \(\operatorname{Gr}^\bullet(X)\) et \(H^{2\bullet}(X,\Zell(\bullet))\)}

\textbf{5.1.} L'exemple donné dans 4.8 n'exclut pas qu'on puisse peut-être trouver, pour un schéma $X$ quelconque, un homomorphisme d'anneaux
\[
 g':\operatorname{Gr}^\bullet(X)\longrightarrow H^{2\bullet}(X,\Zell(\bullet)),
\]
tels que pour tout Module localement libre $E$ sur $X$, $g^i$ transforme la classe de Chern $c_i(E)$ de Exp. V en la classe de Chern cohomologique de SGA 5 VII. Bien entendu, ici $\ell$ désigne un nombre premier distinct des caractéristiques résiduelles de $X$. Comme l'anneau $\operatorname{Gr}^\bullet(X)$ est engendré par les $c^i(E)$, on voit en tous cas que $g'$ est uniquement déterminé par la condition d'être compatible avec les classes de Chern.

\textbf{5.2.} Utilisant l'homomorphisme de Chern
\[
 K^\bullet(X)\longrightarrow\Ch\bigl(H^{2\bullet}(X,\Zell(\bullet))\bigr),
\]
on en conclut, par passage aux gradués associés respectivement à la filtration du premier membre et la filtration par le degré du second, un homomorphisme de groupes gradués (mais pas d'anneaux !)
\[
 \psi:\operatorname{Gr}^\bullet(X)\longrightarrow H^{2\bullet}(X,\Zell(\bullet)),
\]
qui satisfait à la relation
\[
 \psi(c_i(E))=(-1)^{i-1}(i-1)!\,c_i^\ell(E),
\]
où $c_i^\ell(E)$ désigne la classe de Chern $\ell$-adique. On en conclut, par division à partir des $\psi^i$, un homomorphisme d'anneaux
\[
 g_\Q:\operatorname{Gr}^\bullet(X)\longrightarrow
 H^{2\bullet}(X,\Qell(\bullet))=
 H^{2\bullet}(X,\Zell(\bullet))\otimes_{\Zell}\Qell,
\]
transformant classe de Chern en classes de Chern; et tel que l'image du premier membre soit contenue dans celle de $H^2(X,\Zell(\bullet))$; par suite, la question de l'existence de $g'$ se résoud affirmativement pourvu qu'on remplace $H^{2\bullet}(X,\Zell(\bullet))$ par son quotient par son sous-groupe de torsion. Le premier cas douteux pour définir les composantes $g^i$ de l'hypothètique homomorphisme $g$ est celui de $g^3$, on n'arrive à priori qu'à définir $2g^3=\psi^3$. La réponse à la question d'existence de $g^3$ serait négative si par exemple on pouvait trouver trois faisceaux inversibles $L_1,L_2,L_3$ sur $X$ satisfaisant la relation
\[
 (L_1-1)(L_2-1)(L_3-1)=0\quad\text{dans }K^\bullet(X),
\]
%%%%%%%%%%  scan idx 695  |  folio 682  %%%%%%%%%%
mais tels qu'on ait
\[
 c_1(L_1)c_1(L_2)c_1(L_3)\ne0
 \quad\text{dans }H^6(X,\mu_2^{\otimes3}).
\]

\textbf{5.3.} Lorsque $X$ est lisse et quasi-projectif sur un corps $k$, on peut de même se poser la question de l'existence d'un homomorphisme d'anneaux
\[
 g:\operatorname{Gr}^\bullet(X)\longrightarrow A(X)
\]
à valeurs dans l'anneau de Chow $A(X)$, transformant classes de Chern en classes de Chern, et on peut répéter à ce propos les reflexions qui précèdent, qui montrent en particulier comment construire directement $(-1)^{i-1}(i-1)!g^i$ à défaut de $g^i$ lui-même.

\textbf{5.4.} Signalons, pour terminer ce numéro, qu'essentiellement la même question se pose si, au lieu d'un schéma $X$, on considère un espace topologique compact (disons) $X$, ce qui permet encore de considérer le gradué $\operatorname{Gr}^\bullet(X)$ de $K^\bullet(X)$ muni de sa $\lambda$-filtration, et de se proposer de trouver un homomorphisme d'anneaux
\[
 \operatorname{Gr}^\bullet(X)\longrightarrow H^{2\bullet}(X,\Z)
\]
transformant classes de Chern en classes de Chern.

\section*{6. Théorème de Riemann--Roch cohomologique, et homomorphisme de Gysin}

\textbf{6.1.} Pour simplifier, nous nous bornerons ici aux schémas $X$ admettant un Module inversible ample, ce qui implique que l'on n'a pas à distinguer entre le $K'(X)$ naif et l'autre. Commençons par la formule de Riemann--Roch ``avec dénominateurs'', pour un morphisme propre $f:X\to Y$ d'intersection complète. Comme (pour un nombre premier $\ell$ premier aux caractéristiques résiduelles de $Y$) on peut définir les classes de Chern $\ell$-adiques d'un Module localement libre $E$,
\[
 c_i(E)\in H^{2i}(X,\Zell(i))
\]
(SGA 5 VII), correspondant à un homomorphisme
\begin{equation}\tag{6.1}
 K'(X)\longrightarrow \Ch\bigl(H^{2\bullet}(X,\Zell(\bullet))\bigr),
\end{equation}
la question se pose de donner une variante cohomologique de la formule de Riemann--Roch. Le seul ingrédient qui manque pour donner un sens à la formule est un \emph{homomorphisme de Gysin}
%%%%%%%%%%  scan idx 696  |  folio 683  %%%%%%%%%%
\begin{equation}\tag{6.2}
 H^n(X,\Qell(i))\longrightarrow H^{n-2d}(Y,\Qell(i-d)),
\end{equation}
où $d$ est la
dimension relative virtuelle de $d$. Cet homomorphisme étant supposé défini, et tenant compte de la formule de Riemann--Roch sous la forme de Exp. VIII, à valeurs dans $\operatorname{Gr}'(Y)_{\Q}$, on constate que la validité de la formule de Riemann--Roch $\ell$-adique pour le morphisme $f$ équivaut exactement à la commutativité du carré
\begin{equation}\tag{6.3}
\begin{tikzcd}[column sep=large,row sep=large]
 \operatorname{Gr}'(X)_{\Q} \arrow[r] \arrow[d] & H^{2\bullet}(X,\Qell(\bullet)) \arrow[d] \\
 \operatorname{Gr}'(Y)_{\Q} \arrow[r] & H^{2\bullet}(Y,\Qell(\bullet)),
\end{tikzcd}
\end{equation}
où les flèches horizontales se définissent à partir de (6.1) en passant aux gradués associés et en divisant, de façon à transformer classes de Chern en classes de Chern, et où les flèches verticales sont respectivement la flèche de Exp. VIII et l'homomorphisme de Gysin en cohomologie (supposé donné).

\textbf{6.2.} Notons que l'homomorphisme de Gysin est défini en tous cas si $Y$ est lisse sur un corps $k$ (SGA 5 IV), en utilisant le théorème de pureté cohomologique relatif sur $k$; la même construction s'applique lorsque $Y$ est excellent régulier et de caractéristique nulle, grâce au théorème de pureté cohomologique absolu de SGA 4 XIX, et elle s'appliquera automatiquement pour tout $Y$ excellent et régulier, sans restriction de caractéristique, une fois le théorème de pureté cohomologique absolu prouvé en toute généralité. Enfin, la même construction s'applique, grâce au théorème de pureté cohomologique relatif (SGA 4 XVI 3.7), lorsque $X$ et $Y$ sont tous deux lisses sur un même schéma $S$ et que $f$ est un $S$-morphisme. Il y a lieu de conjecturer que dans chacun de ces cas où l'homomorphisme de Gysin est défini, le diagramme (6.3) est bien commutatif. Le seul cas où cette assertion soit actuellement vérifiée est celui où $X$ et $Y$ sont tous les deux lisses sur un corps $k$, et résulte alors de la commutativité du diagramme analogue à (6.3)
\begin{equation}\tag{6.4}
\begin{tikzcd}[column sep=large,row sep=large]
 A(X) \arrow[r] \arrow[d] & H^{2\bullet}(X,\Zell(\bullet)) \arrow[d] \\
 A(Y) \arrow[r] & H^{2\bullet}(Y,\Zell(\bullet)),
\end{tikzcd}
\end{equation}
%%%%%%%%%%  scan idx 697  |  folio 684  %%%%%%%%%%
faisant intervenir les anneaux de classes de cycles (mod. équivalence rationelle) $A(X)$ et $A(Y)$, commutativité établie dans SGA 5 IV.

\textbf{6.3.} Notons que dans chacun des cas envisagés plus haut où on arrive à définir un homomorphisme de Gysin, on obtient en fait une construction pour la cohomologie à coefficients dans $\Zell(\bullet)$, et non seulement dans $\Qell(\bullet)$. Cela permet, lorsque $f$ est une \emph{immersion}, de poser la question de la validité de la formule de Riemann--Roch ``sans dénominateurs'', comme il est expliqué au no. 3. Lorsque $X$ et $Y$ sont tous deux lisses sur un corps $k$, la commutativité de (6.4) nous montre qu'il suffirait de prouver la formule sans dénominateurs dans l'anneau de Chow $A(Y)$ cf. (3.3).

\textbf{6.4.} Lorsqu'on ne fait pas d'hypothèse de régularité sur $Y$, et même si $X$ et $Y$ sont de type fini sur le corps $\C$, on ne voit plus de définition plausible d'un homomorphisme de Gysin cohomologique (6.2), même pour une immersion régulière de codimension $1$. Peut-on en trouver un néanmoins avec des propriétés raisonnables, en particulier rendant commutatif le carré (6.3)? La question est liée évidemment à celle de définir, pour un $X$ d'intersection complète de codimension $i$ dans $Y$, ``sa classe de cohomologie''
\[
 \gamma(X)\in H^{2i}(Y,\Zell(i)),
\]
de façon à tranformer intersections (sans composantes excédentaires) en cup-produits. Cette dernière question se résoud affirmativement tout au moins modulo torsion, i.e. dans $H^{2\bullet}(Y,\Qell(\bullet))$, en définissant
\[
 \gamma(X)=1/(-1)^{i-1}(i-1)!\,c_i(\mathcal O_X).
\]
Mais j'ignore si $c_i(\mathcal O_X)$ est divisible par $(i-1)!$ dans $H^{2i}(Y,\Zell(i))$. En particulier, si $Y$ est un schéma projectif de dimension $3$ sur le corps des complexes, disons, et $X$ réduit à un point (singulier) de $Y$, $c_3(\mathcal O_X)$ est-il divisible par $2$ dans $H^6(Y,\Z_2(3))$, et en particulier la classe de Chern est elle nulle dans $H^6(Y,\Z/2\Z)$? Signalons également dans la même voie la question suivante: soit $Y$ un schéma projectif (singulier) de dimension $4$, $X$ un sous-schéma fermé de dimension $1$ qui soit globalement une intersection de trois diviseurs positifs $D_1,D_2,D_3$, considérons le produit de leurs classes de Chern mod. $2$:
\[
 c_1(D_1)c_1(D_2)c_1(D_3)\in H^6(Y,\Z/2\Z);
\]
%%%%%%%%%%  scan idx 698  |  folio 685  %%%%%%%%%%
cet
élément du $H^6$ est-il indépendant du choix particulier des $D_i$ d'intersection $X$?

\section*{7. Classes de Chern des complexes parfaits}

\textbf{7.1.} Lorsqu'on se pose la question de développer une formule de Riemann--Roch cohomologique, avec ou sans dénominateurs, en dehors de l'hypothèse assez artificielle d'existence d'un Module inversible ample sur les schémas envisagés, on est nécessairement confronté à la question d'une définition, pour un complexe de Modules parfait $L^\bullet$ sur un schéma $X$, des classes de Chern
\[
 c_i(L^\bullet)\in H^{2i}(X,\Zell(i)),
\]
généralisant les classes de Chern des Modules localement libres SGA 5 VII. La même question se pose d'ailleurs pour des topos localement annelés généraux (comparer [10]), ainsi que pour les complexes parfaits sur des espaces paracompacts annelés par les fonctions complexes continues, en travaillant alors avec la cohomologie entière ordinaire. Dans le cas d'un espace compact, on sait que le $K'(X)$ construit avec les complexes parfaits est le même que le ``naif'' construit avec les Modules localement libres (Exp. IV), et dans ce cas on est ramené trivialement à la définition classique. Mais déjà dans le cas d'un espace localement compact le problème posé semble non trivial, puisque la cohomologie entière d'un tel espace n'est pas toujours la limite projective des cohomologies de ses parties compactes.

\textbf{7.2.} Plus précisément, lorsque $Y$ est une partie fermée de $X$ telle que $L^\bullet|(X-Y)$ soit nul (au sens des catégories dérivées), i.e. telle que les faisceaux de cohomologie de $L^\bullet$ soient à supports dans $Y$, on aimerait définir des classes de Chern \emph{à support dans} $Y$
\[
 c_i^Y(L^\bullet)\in H_Y^{2i}(X,\Zell(i))\quad\text{resp.}\quad H_Y^{2i}(X,\Z),
\]
qui redonneraient les classes de Chern ``ordinaires'' (sic) de 7.1 à l'aide de l'homomorphisme canonique
\[
 H_Y^\bullet(X,-)\longrightarrow H^\bullet(X,-).
\]
%%%%%%%%%%  scan idx 699  |  folio 686  %%%%%%%%%%
Même
dans le cas où ces dernières se ramènent aux classes de Chern des Modules localement libres ($X$ schéma avec Module inversible ample, ou espace compact), la construction de ces classes de Chern ``à supports'' n'est pas faite, et ne pourra sans doute se faire qu'en même temps que celle des classes de Chern générales postulées dans 7.1.\footnote{Ajouté en octobre 1968. Les problèmes soulevés ici ont été résolus dans le cadre de la \emph{cohomologie de Hodge} par L. Illusie; cf. [13].}

Il est probable que les classes de Chern avec supports auront un rôle à jouer dans la démonstration de la formule de Riemann--Roch cohomologique sans dénominateurs (cf. 3.5). D'autre part, elles devraient permettre de donner une version améliorée de la formule de Riemann--Roch (avec ou sans dénominateurs) en introduisant dans cette formule des familles de supports.

\section*{8. Anneau de Chow des schémas réguliers}

Le problème est de définir l'équivalence rationelle des cycles, et d'étudier l'anneau de classes obtenu (notamment ses propriétés fonctorielles), pour un schéma noethérien régulier $X$ admettant un Module inversible ample, en généralisant la théorie connue lorsque $X$ est lisse sur un corps [4]. Le résultat technique clef sera encore un ``moving lemma'', assurant qu'on peut faire varier un cycle dans sa classe de façon à intersecter avec la bonne dimension un cycle donné.

Le problème envisagé est assez naturel à priori, d'autant plus que nous avons vu (4.5 à 4.8) que les substituts $\operatorname{Gr}'(X)$, $\Gtop^\bullet(X)$ de l'anneau de Chow $A(X)$ ne sont guère satisfaisants quand on ne se borne pas à une théorie ``modulo torsion''. Lorsque $X$ se déduit d'un schéma algébrique projectif et lisse $Y$ sur un corps $k$, en prenant le schéma local au sommet d'un cône projetant $C$ et enlevant l'origine, on trouve (en admettant que $A(X)$ est bien la limite inductive des anneaux de Chow des $U-(O)$, où $U$ parcourt les voisinages ouverts dans $C$ du sommet $O$) que $A(X)$ est canoniquement isomorphe à l'anneau quotient $A(Y)/\xi A(Y)$, où $\xi\in A^1(Y)$ est induit par l'immersion projective considérée de $Y$. Cela montre que la connaissance de l'anneau de Chow $A(X)$ (de type non classique) équivaut à la connaissance de $A(Y)/\xi A(Y)$, qu'on pourrait appeler la ``partie primitive'' de l'anneau de Chow $A(Y)$ (en un sens suggéré par la théorie cohomologique de Lefschetz--Hodge des variétés algébriques projectives non singulières). Cet exemple montre que les ``anneaux de Chow locaux'' des anneaux locaux à singularités
%%%%%%%%%%  scan idx 700  |  folio 687  %%%%%%%%%%
isolées sont des invariants géométriques tout à fait non triviaux, et qui auront sans doute leur rôle à jouer dans l'étude ``arithmétiques'' de ces anneaux locaux.

Bien entendu, une fois en possession d'une notion raisonnable d'anneau de Chow pour un schéma noethérien régulier admettant un Module inversible ample, il y a lieu de généraliser à ce cas les relations établies resp. conjecturées au no. 4 entre cet anneau et l'anneau $\Gtop^\bullet(X)$, et à formuler en termes de cet anneau la question de la ``formule de Riemann--Roch sans dénominateurs'' énoncée au no. 3.

\subsection*{Bibliographie de l'Exposé XIV}
\begin{enumerate}[label={[\arabic*]}]
\item M. F. Atiyah et F. Hirzebruch, \emph{The Riemann--Roch theorem for analytic embeddings}, Topology, vol. 1, pp. 151--166 (1962).
\item A. Borel et J.-P. Serre, \emph{Le théorème de Riemann--Roch}, Bull. Soc. Math. France, vol. 86, pp. 97--136 (1958).
\item R. Bott et H. Samelson, Amer. Journ. of Math. vol. 80, p. 1004 (1958).
\item C. Chevalley, \emph{Les classes d'équivalence rationelle}, I et II, Séminaire à l'ENS 1958, exposés 2 et 3.
\item C. Chevalley, \emph{Classification des groupes de Lie algébriques}, Séminaire à l'ENS 1956/58.
\item A. Dold et D. Puppe, \emph{Homologie nicht-additiver Funktoren. Anwendungen}, Annales de l'Inst. Fourier, t. XI, pp. 201--312 (1961).
\item H. Grauert, \emph{Ein Theorem der analytischen Garbentheorie}, Pub. Math. no. 5, pp. 5--64 (1960).
\item A. Grothendieck, \emph{Classes de faisceaux et théorème de Riemann--Roch}, notes multigraphiées par les soins de l'Institute for Advanced Study, 1957 (travail cité [RRR] dans le présent exposé, et qui est reproduit dans notre séminaire sous forme d'Appendice au Chap. O).
\item A. Grothendieck, \emph{La théorie des classes de Chern}, Bull. Soc. Math. France, vol. 86, pp. 137--154 (1958).
\item A.
%%%%%%%%%%  scan idx 701  |  folio 688  %%%%%%%%%%
Grothendieck, \emph{Classes de Chern et représentations linéaires des groupes discrets}, in : \emph{Dix exposés sur la cohomologie des schémas}, Masson--North Holland Pub. Cie.
\item H. Hironaka, \emph{Resolution of singularities of an algebraic variety over a field of characteristic zero}, Ann. of Math. 79, pp. 109--326 (1964).
\item L. Hodgkin, thèse (à paraitre ?).
\item L. Illusie, travail en préparation.
\item J.-P. Serre, \emph{Algébre Locale, Multiplicités}, Lecture Notes in Math. no. 11 (1965) Springer.
\item A. Borel, \emph{Sous-groupes commutatifs et torsion des groupes de Lie compacts connexes}, Tohoku Math. Journal, vol. 13, pp. 216--240 (1961).
\item[{[16]}(*)] P. Deligne, Exp. XVIII in \emph{Théorie des topos et cohomologie étale des schémas (SGA 6)}, par M. Artin, A. Grothendieck et J. L. Verdier.
\item[{[17]}(*)] Yu. I. Manin, \emph{Correspondances, motifs et transformations monoïdales} (en russe). Mat. Sbornik, t. 77 (119), n\textsuperscript{o} 4, 1968.
\end{enumerate}

\noindent{\footnotesize (*) Ajouté en février 1969.}

%%%%%%%%%%  scan idx 702  |  folio 689  %%%%%%%%%%
\subsection*{Note pour la page 1}

La théorie des opérations $\Lambda^i$ dans $D^-(X)$, et des opérations $\lambda^i$ dans $K'(X)$, a été écrite par P. Deligne en 1968. Il prouve que $K'(X)$ devient ainsi un pré-$\lambda$-anneau, et lorsque $\mathcal O_X$ est un faisceau d'Algèbres sur un corps, que $K'(X)$ est même un $\lambda$-anneau. Il est plausible qu'un raffinement de la méthode de Deligne permette de se débarasser de l'hypothèse restrictive sur $\mathcal O_X$.

\subsection*{Notes pour la page 11}

\noindent(*) (Avril 1969) D. MUMFORD nous signale une démonstration de (9.6) dans l'\emph{anneau de Chow} (sans négliger la torsion), qu'il a obtenue dans un séminaire à Harvard en 1959 ! Le lecteur trouvera la démonstration de ce beau résultat dans SGA 5 VII 9. Ce résultat contient évidemment le résultat cohomologique correspondant, signalé dans la note (**). Cependant, le résultat cohomologique est prouvé par J. P. JOUANOLOU sous des conditions plus générales (sans hyperplan quasi-projectif, et dans le cas lisse relatif sur une base quelconque).

\medskip

\noindent(**) (janvier 1969) J. P. JOUANOLOU vient de prouver ces formules pour tout $y$ sans négliger la torsion, cf. SGA 5 VII \S{} 4.9. DELIGNE et MANIN 17 avaient auparavant remarqué indépendamment, et par des arguments (faciles) de nature assez différente, qu'il suffisait d'établir des formules pour $y$ ``algébrique'', et même pour $y=1$ dans 16, pour conclure la même relation dans le cas général.



% Final-index helper macros
\providecommand{\Qcoh}{\operatorname{Qcoh}}
\providecommand{\Coh}{\operatorname{Coh}}
\providecommand{\Parf}{\operatorname{Parf}}
\providecommand{\Lib}{\operatorname{Lib}}
\providecommand{\Loclib}{\operatorname{Loclib}}
